% !TeX program = xelatex
% Font dan ukuran kertas mengikuti PDF EAS asli (8,5 x 13 inci).
\documentclass[12pt]{article}
\usepackage{amsmath,amsthm}
\usepackage{fontspec}
\usepackage{unicode-math}
\setmainfont{Times New Roman}
\setmathfont{Cambria Math}
\usepackage{graphicx}
\graphicspath{{../../../PPT Kalkulus/Gambar & Logo/}}
\usepackage{array,tabularx,enumitem}
\usepackage{fancyhdr}
\usepackage[most]{tcolorbox}
\usepackage{hyperref}
\hypersetup{colorlinks=true,linkcolor=blue,urlcolor=cyan,
    pdftitle={Bahas Soal EAS Kalkulus 1 2025/2026},
    pdfpagemode=UseOutlines,hypertexnames=false}
\usepackage{geometry}
\geometry{paperwidth=612bp,paperheight=936bp,
    left=54bp,right=54bp,top=143bp,bottom=55bp,
    headheight=120bp,headsep=8bp,footskip=25bp}
\setlength{\parindent}{0pt}
\setlength{\parskip}{0pt}
\setlength{\tabcolsep}{5bp}
\setlength{\jot}{3bp}
\setlist{topsep=4bp,itemsep=3bp,parsep=0pt,partopsep=0pt}
\setlist[enumerate,1]{label=\arabic*.,leftmargin=29bp,labelsep=8bp}
\setlist[enumerate,2]{label=\alph*.,leftmargin=18bp,labelsep=5bp}
\setlist[itemize]{leftmargin=18bp}
\allowdisplaybreaks[1]
\newcommand{\R}{\mathbb{R}}
\newcommand{\N}{\mathbb{N}}
\newcommand{\C}{\mathbb{C}}
\newcommand{\Z}{\mathbb{Z}}
\newcommand{\Q}{\mathbb{Q}}
\newcommand{\cis}{\text{cis}}

\pagestyle{fancy}
\fancyhf{}
\renewcommand{\headrulewidth}{0.4bp}
\renewcommand{\footrulewidth}{0.4bp}
\newcommand{\packetid}{1}
\newcommand{\answername}{Jawaban}
\newcommand{\startpacket}[1]{%
    \clearpage
    \renewcommand{\packetid}{#1}%
    \setcounter{page}{1}%
}

% Kepala ujian diulang pada setiap halaman, termasuk lanjutan jawaban.
\newcommand{\examhead}[7]{%
    \begin{minipage}[b]{124bp}
        \includegraphics[width=117bp]{logoITSBiru.png}\par\vspace{8bp}
        {\fontsize{20bp}{22bp}\selectfont\bfseries SKPB -- ITS}
    \end{minipage}\hfill
    \begin{minipage}[b]{370bp}
        {\fontsize{16bp}{17bp}\selectfont\bfseries\underline{#1}}\par
        {\fontsize{12bp}{14bp}\selectfont
        \begin{tabular}[b]{@{}l@{\hspace{12bp}:\hspace{5bp}}l@{}}
            #2 & #3\\
            #4\\
            #5\\
            #6\\[3bp]
            #7
        \end{tabular}}
    \end{minipage}%
}
\newcommand{\examheaderID}[4]{%
    \renewcommand{\answername}{Jawaban}%
    \fancyhead[C]{\examhead
        {EVALUASI AKHIR SEMESTER GANJIL 2025/2026}
        {Mata Kuliah, SKS}{#1}
        {Hari, Tanggal & #2}
        {Jam, Waktu Ujian & #3}
        {Kelas, Sifat & #4}
        {{\fontsize{16bp}{20bp}\selectfont\bfseries Nama} & \\[6bp]
         {\fontsize{16bp}{20bp}\selectfont\bfseries NRP, Kelas} &}}%
    \fancyfoot[R]{\fontsize{12bp}{14bp}\selectfont Halaman \thepage\ dari \pageref*{akhir-paket-\packetid}}%
}
\newcommand{\examheaderEN}[4]{%
    \renewcommand{\answername}{Solution}%
    \fancyhead[C]{\examhead
        {FINAL EXAM OF THE ODD SEMESTER 2025/2026}
        {Course, Credit}{#1}
        {Day, Date & #2}
        {Time, Exam Time & #3}
        {Class, Type & #4}
        {{\fontsize{16bp}{20bp}\selectfont\bfseries Name} & \\[6bp]
         {\fontsize{16bp}{20bp}\selectfont\bfseries NRP, Class} &}}%
    \fancyfoot[R]{\fontsize{12bp}{14bp}\selectfont Page \thepage\ of \pageref*{akhir-paket-\packetid}}%
}

\newcommand{\warningsID}{%
    \vspace*{-8bp}%
    {\centering\fontsize{12bp}{14bp}\selectfont
    Tuliskan Nama, NRP, dan Nomor Kelas pada setiap lembar kerja ini\\
    \textbf{DILARANG MEMBAWA/MENGGUNAKAN KALKULATOR DAN ALAT KOMUNIKASI}\\
    \textbf{DILARANG MEMBERIKAN/MENERIMA JAWABAN SELAMA UJIAN BERLANGSUNG}\\
    ``Setiap tindak kecurangan akan mendapatkan sanksi akademik''\par}
    \vspace{2bp}\hrule height 0.4bp\vspace{4bp}%
}
\newcommand{\warningsEN}{%
    \vspace*{-8bp}%
    {\centering\fontsize{12bp}{14bp}\selectfont
    Write your Name, Student ID (NRP), and Class Number on every sheet of this worksheet\\
    {\fontsize{11bp}{14bp}\selectfont USING OR BRINGING CALCULATORS AND COMMUNICATION DEVICES IS STRICTLY PROHIBITED}\\
    SHARING OR RECEIVING ANSWERS DURING THE EXAM IS STRICTLY PROHIBITED\\
    ``Any act of dishonesty will result in academic sanctions''\par}
    \vspace{2bp}\hrule height 0.4bp\vspace{4bp}%
}
\newcommand{\cpmktableID}{%
    {\renewcommand{\arraystretch}{1}\fontsize{12bp}{14bp}\selectfont\hyphenpenalty=10000
    \begin{tabularx}{\textwidth}{|p{48bp}|X|>{\centering\arraybackslash}p{41bp}|>{\centering\arraybackslash}p{56bp}|}
        \hline
        CPMK & Kemampuan & No Soal & Bobot (\%)\\\hline
        CPMK-2 & Mampu menentukan kekontinuan fungsi dan turunannya & 2 & 20\\\hline
        CPMK-3 & Mampu menghitung integral dan memahami Teorema Fundamental Kalkulus & 1, 3 & 40\\\hline
        CPMK-4 & Mampu mengaplikasikan bentuk peubah kompleks dalam bentuk polar serta mencari akar-akar persamaannya & 4 & 20\\\hline
        CPMK-5 & Mampu menerapkan konsep matriks untuk menyelesaikan sistem persamaan linier dan menentukan nilai eigen & 5 & 20\\\hline
    \end{tabularx}\par}\vspace{4bp}%
}
\newcommand{\cpmktableEN}{%
    {\renewcommand{\arraystretch}{1}\fontsize{12bp}{14bp}\selectfont\hyphenpenalty=10000
    \begin{tabularx}{\textwidth}{|p{48bp}|X|>{\centering\arraybackslash}p{50bp}|>{\centering\arraybackslash}p{63bp}|}
        \hline
        CLO & Ability & Question & Weight (\%)\\\hline
        CLO-2 & Able to determine the continuity of functions and their derivatives & 2 & 20\\\hline
        CLO-3 & Able to calculate integrals and understand the Fundamental Theorem of Calculus & 1, 3 & 40\\\hline
        CLO-4 & Able to apply complex variables in polar form and find the roots of the equations & 4 & 20\\\hline
        CLO-5 & Able to apply matrix concepts to solve systems of linear equations and determine eigenvalues & 5 & 20\\\hline
    \end{tabularx}\par}\vspace{4bp}%
}
\newcommand{\scoretable}[2]{%
    {\renewcommand{\arraystretch}{1}\fontsize{12bp}{14bp}\selectfont
    \begin{tabularx}{\textwidth}{|>{\centering\arraybackslash}p{60bp}|*{5}{>{\centering\arraybackslash}X|}>{\centering\arraybackslash}p{60bp}|}
        \hline
        #1 & 1 & 2 & 3 & 4 & 5 & TOTAL\\\hline
        #2 & & & & & & \rule[-10bp]{0pt}{25bp}\\\hline
    \end{tabularx}\par}\vspace{4bp}%
}
\newcommand{\scoretableID}{\scoretable{Nomor}{Nilai}}
\newcommand{\scoretableEN}{\scoretable{Question}{Score}}

% Satu soal beserta jawaban dalam satu bingkai; jawaban panjang dapat bersambung.
\newenvironment{soaljawab}[1]{%
    \ifnum#1>1\clearpage\fi
    \begin{tcolorbox}[enhanced,breakable,height fill,height fixed for=all,
        sharp corners,colback=white,colframe=black,boxrule=0.4bp,
        boxsep=0pt,left=12bp,right=12bp,top=6bp,bottom=10bp,
        before skip=0pt,after skip=0pt,pad at break=8bp]
    \fontsize{12bp}{14bp}\selectfont
    \setlength{\abovedisplayskip}{8bp}
    \setlength{\belowdisplayskip}{8bp}
    \setlength{\abovedisplayshortskip}{8bp}
    \setlength{\belowdisplayshortskip}{8bp}
    \renewcommand{\arraystretch}{1.15}
    \begin{enumerate}[start=#1,topsep=0pt]
}{%
    \end{enumerate}
    \end{tcolorbox}
}
\newcommand{\jawaban}{\par\medskip\textbf{\answername.}\par\smallskip}

\begin{document}


% PAKET 1: Senin, 8 Desember 2025 | 07.00 - 08.40 WIB | Reguler (Kelas 40-49, 62-66)
\startpacket{1}
\pdfbookmark[0]{Paket 1: Senin 07.00 Reguler (Kelas 40--49, 62--66)}{paket1}
\examheaderID{Kalkulus 1, 3 SKS}{Senin, 8 Desember 2025}{07:00 -- 08:40, 100 menit}{40 -- 49 dan 62 -- 66, Tertutup}
\warningsID
\cpmktableID
\scoretableID
\begin{soaljawab}{1}
\item Hitunglah integral berikut
        \[
          \int \cos(3x)\sqrt{2 - \sin(3x)} \, dx.
        \]
\jawaban
Gunakan metode substitusi. Misalkan $u = 2 - \sin(3x)$, maka:
        \[
          du = -3\cos(3x)\,dx \implies \cos(3x)\,dx = -\frac{1}{3}\,du.
        \]
        Sehingga integralnya menjadi:
        \begin{align*}
          \int \cos(3x)\sqrt{2 - \sin(3x)} \, dx & = \int \sqrt{u} \left(-\frac{1}{3}\right) du = -\frac{1}{3} \int u^{1/2} \, du \\
                                                 & = -\frac{1}{3} \cdot \frac{2}{3} u^{3/2} + C                                   \\
                                                 & = -\frac{2}{9} (2 - \sin(3x))^{3/2} + C.
        \end{align*}
\end{soaljawab}

\begin{soaljawab}{2}
\item Carilah maksimum dan minimum absolut dari
        \[
          y = \frac{2x + 5}{x^2 - 4}
        \]
        pada interval $[-5, -3]$.
\jawaban
Diberikan fungsi $y = f(x) = \dfrac{2x + 5}{x^2 - 4}$ pada $[-5, -3]$.
        Perhatikan bahwa penyebut bernilai nol saat $x = \pm 2$, yang berada di luar interval $[-5, -3]$. Jadi $f(x)$ kontinu dan terdiferensialkan pada $[-5, -3]$.

        Cari turunan pertama:
        \begin{align*}
          f'(x) & = \frac{2(x^2 - 4) - (2x + 5)(2x)}{(x^2 - 4)^2} = \frac{2x^2 - 8 - 4x^2 - 10x}{(x^2 - 4)^2}                          \\
                & = \frac{-2x^2 - 10x - 8}{(x^2 - 4)^2} = \frac{-2(x^2 + 5x + 4)}{(x^2 - 4)^2} = \frac{-2(x + 1)(x + 4)}{(x^2 - 4)^2}.
        \end{align*}
        Titik kritis dari $f'(x) = 0$ adalah $x = -1$ dan $x = -4$.
        Titik yang berada di dalam interval $[-5, -3]$ hanyalah $x = -4$.

        Evaluasi nilai fungsi pada titik kritis dan titik ujung interval:
        \begin{itemize}
          \item $x = -5$: $f(-5) = \dfrac{2(-5) + 5}{(-5)^2 - 4} = \dfrac{-5}{21} \approx -0.238$
          \item $x = -4$: $f(-4) = \dfrac{2(-4) + 5}{(-4)^2 - 4} = \dfrac{-3}{12} = -\frac{1}{4} = -0.25$
          \item $x = -3$: $f(-3) = \dfrac{2(-3) + 5}{(-3)^2 - 4} = \dfrac{-1}{5} = -0.2$
        \end{itemize}
        Dari nilai-nilai di atas:
        \begin{itemize}
          \item Nilai maksimum absolut adalah $-\dfrac{1}{5}$ pada $x = -3$.
          \item Nilai minimum absolut adalah $-\dfrac{1}{4}$ pada $x = -4$.
        \end{itemize}
\end{soaljawab}

\begin{soaljawab}{3}
\item Gambarkan grafik $y = |2x + 1|$ untuk $-1 \le x \le 4$ dan hitung integral
        \[
          \int_{-1}^{4} |2x + 1| \, dx
        \]
        dengan menggunakan rumus dari geometri bidang yang sesuai.
\jawaban
Fungsi $y = |2x + 1|$ bernilai:
        \[
          |2x + 1| = \begin{cases}
            2x + 1,    & x \ge -\frac{1}{2} \\
            -(2x + 1), & x < -\frac{1}{2}
          \end{cases}
        \]
        Pada interval $[-1, 4]$, daerah di bawah kurva terbagi menjadi dua segitiga:
        \begin{enumerate}
          \item Segitiga I pada $\left[-1, -\frac{1}{2}\right]$:
                \begin{itemize}
                  \item Alas $a_1 = -\frac{1}{2} - (-1) = \frac{1}{2}$
                  \item Tinggi $t_1 = |2(-1) + 1| = 1$
                  \item Luas $L_1 = \frac{1}{2} \cdot \frac{1}{2} \cdot 1 = \frac{1}{4}$
                \end{itemize}
          \item Segitiga II pada $\left[-\frac{1}{2}, 4\right]$:
                \begin{itemize}
                  \item Alas $a_2 = 4 - \left(-\frac{1}{2}\right) = \frac{9}{2}$
                  \item Tinggi $t_2 = |2(4) + 1| = 9$
                  \item Luas $L_2 = \frac{1}{2} \cdot \frac{9}{2} \cdot 9 = \frac{81}{4}$
                \end{itemize}
        \end{enumerate}
        Maka dengan rumus geometri bidang:
        \[
          \int_{-1}^{4} |2x + 1| \, dx = L_1 + L_2 = \frac{1}{4} + \frac{81}{4} = \frac{82}{4} = \frac{41}{2} = 20.5.
        \]
\end{soaljawab}

\begin{soaljawab}{4}
\item Nyatakan bilangan kompleks $\sqrt{3} + i$ dalam bentuk kutub. Kemudian carilah semua bilangan kompleks $z$ yang memenuhi
        \[
          z = (\sqrt{3} + i)^{1/3}.
        \]
\jawaban
Misalkan $w = \sqrt{3} + i$.
        Modulus $r = |w| = \sqrt{(\sqrt{3})^2 + 1^2} = \sqrt{3 + 1} = 2$.
        Argumen $\theta = \arctan\left(\frac{1}{\sqrt{3}}\right) = 30^\circ = \frac{\pi}{6}$ (karena berada di kuadran I).
        Sehingga bentuk kutubnya:
        \[
          \sqrt{3} + i = 2 \operatorname{cis}(30^\circ) = 2 \left(\cos\frac{\pi}{6} + i\sin\frac{\pi}{6}\right).
        \]
        Akar-akar $z = w^{1/3}$ untuk $k = 0, 1, 2$:
        \[
          z_k = \sqrt[3]{2} \operatorname{cis}\left(\frac{30^\circ + 360^\circ k}{3}\right) = 2^{1/3} \operatorname{cis}(10^\circ + 120^\circ k).
        \]
        Diperoleh ketiga akar:
        \begin{align*}
          z_0 & = 2^{1/3} \operatorname{cis}(10^\circ),  \\
          z_1 & = 2^{1/3} \operatorname{cis}(130^\circ), \\
          z_2 & = 2^{1/3} \operatorname{cis}(250^\circ).
        \end{align*}
\end{soaljawab}

\begin{soaljawab}{5}
\item Carilah invers dari
        \[
          A = \begin{bmatrix}
            1  & -2 & -6 \\
            2  & -1 & 4  \\
            -2 & 2  & 1
          \end{bmatrix}
        \]
        dengan menggunakan operasi baris elementer.
\jawaban
Gabungkan $A$ dengan matriks identitas $I_3$. Lakukan OBE pada
kedua bagian hingga $[A\mid I_3]$ menjadi $[I_3\mid A^{-1}]$.
Di sini $B_i$ menyatakan baris ke-$i$; label $B_i+kB_j$ berarti
mengganti $B_i$ dengan $B_i+kB_j$.
{\setlength{\arraycolsep}{4bp}
\begin{align*}
    & \left[\begin{array}{ccc|ccc}
            1 & -2 & -6 & 1 & 0 & 0 \\
            2 & -1 & 4 & 0 & 1 & 0 \\
            -2 & 2 & 1 & 0 & 0 & 1
        \end{array}\right]
        \begin{array}{c}\xrightarrow{B_2-2B_1}\end{array}
        \left[\begin{array}{ccc|ccc}
            1 & -2 & -6 & 1 & 0 & 0 \\
            0 & 3 & 16 & -2 & 1 & 0 \\
            -2 & 2 & 1 & 0 & 0 & 1
        \end{array}\right] \\[6bp]
    & \begin{array}{c}\xrightarrow{B_3+2B_1}\end{array}
        \left[\begin{array}{ccc|ccc}
            1 & -2 & -6 & 1 & 0 & 0 \\
            0 & 3 & 16 & -2 & 1 & 0 \\
            0 & -2 & -11 & 2 & 0 & 1
        \end{array}\right]
        \begin{array}{c}\xrightarrow{B_2+B_3}\end{array}
        \left[\begin{array}{ccc|ccc}
            1 & -2 & -6 & 1 & 0 & 0 \\
            0 & 1 & 5 & 0 & 1 & 1 \\
            0 & -2 & -11 & 2 & 0 & 1
        \end{array}\right] \\[6bp]
    & \begin{array}{c}\xrightarrow{B_3+2B_2}\end{array}
        \left[\begin{array}{ccc|ccc}
            1 & -2 & -6 & 1 & 0 & 0 \\
            0 & 1 & 5 & 0 & 1 & 1 \\
            0 & 0 & -1 & 2 & 2 & 3
        \end{array}\right]
        \begin{array}{c}\xrightarrow{B_1+2B_2}\end{array}
        \left[\begin{array}{ccc|ccc}
            1 & 0 & 4 & 1 & 2 & 2 \\
            0 & 1 & 5 & 0 & 1 & 1 \\
            0 & 0 & -1 & 2 & 2 & 3
        \end{array}\right] \\[6bp]
    & \begin{array}{c}\xrightarrow{-B_3}\end{array}
        \left[\begin{array}{ccc|ccc}
            1 & 0 & 4 & 1 & 2 & 2 \\
            0 & 1 & 5 & 0 & 1 & 1 \\
            0 & 0 & 1 & -2 & -2 & -3
        \end{array}\right]
        \begin{array}{c}\xrightarrow{B_1-4B_3}\end{array}
        \left[\begin{array}{ccc|ccc}
            1 & 0 & 0 & 9 & 10 & 14 \\
            0 & 1 & 5 & 0 & 1 & 1 \\
            0 & 0 & 1 & -2 & -2 & -3
        \end{array}\right] \\[6bp]
    & \begin{array}{c}\xrightarrow{B_2-5B_3}\end{array}
        \left[\begin{array}{ccc|ccc}
            1 & 0 & 0 & 9 & 10 & 14 \\
            0 & 1 & 0 & 10 & 11 & 16 \\
            0 & 0 & 1 & -2 & -2 & -3
        \end{array}\right]
\end{align*}}
Bagian kiri sudah menjadi $I_3$, sehingga bagian kanan merupakan inversnya:
\[
    A^{-1} = \begin{bmatrix}
            9 & 10 & 14 \\
            10 & 11 & 16 \\
            -2 & -2 & -3
        \end{bmatrix}.
\]
\end{soaljawab}

\label{akhir-paket-1}

% PAKET 2: Senin, 8 Desember 2025 | 09.00 - 10.40 WIB | IUP (Class 1, 4, 67, 70, 71)
\startpacket{2}
\pdfbookmark[0]{Paket 2: Monday 09.00 IUP (Class 1, 4, 67, 70, 71)}{paket2}
\examheaderEN{Kalkulus 1, 3 SKS}{Monday, 8 December 2025}{09:00 -- 10:40, 100 minutes}{1, 4, 67, 70, 71, Closed book}
\warningsEN
\cpmktableEN
\scoretableEN
\begin{soaljawab}{1}
\item Calculate the following integral
        \[
          \int \frac{\cos(3x)}{\sin^6(3x)} \, dx.
        \]
\jawaban
Let $u = \sin(3x)$, then $du = 3\cos(3x)\,dx \implies \cos(3x)\,dx = \frac{1}{3}\,du$.
        \begin{align*}
          \int \frac{\cos(3x)}{\sin^6(3x)} \, dx & = \int u^{-6} \left(\frac{1}{3}\right) du = \frac{1}{3} \cdot \left(\frac{u^{-5}}{-5}\right) + C \\
                                                 & = -\frac{1}{15\sin^5(3x)} + C.
        \end{align*}
\end{soaljawab}

\begin{soaljawab}{2}
\item A particle moves east at a speed of $4\text{ m/s}$, while another particle moves north at a speed of $3\text{ m/s}$. At a certain moment, the first particle is $8\text{ m}$ from the origin, and the second particle is $6\text{ m}$ from the origin. Determine the rate of change of the distance between the two particles at that moment.
\jawaban
Let the position of the first particle moving east be $x(t)$, and the second moving north be $y(t)$.
        Given:
        \[
          \frac{dx}{dt} = 4\text{ m/s}, \quad \frac{dy}{dt} = 3\text{ m/s}.
        \]
        At the given instant, $x = 8\text{ m}$ and $y = 6\text{ m}$.
        The distance $D$ between the particles satisfies:
        \[
          D^2 = x^2 + y^2 \implies D = \sqrt{8^2 + 6^2} = 10\text{ m}.
        \]
        Differentiating implicitly with respect to time $t$:
        \[
          2D \frac{dD}{dt} = 2x \frac{dx}{dt} + 2y \frac{dy}{dt} \implies D \frac{dD}{dt} = x \frac{dx}{dt} + y \frac{dy}{dt}.
        \]
        Substituting the known values:
        \[
          10 \frac{dD}{dt} = 8(4) + 6(3) = 32 + 18 = 50 \implies \frac{dD}{dt} = \frac{50}{10} = 5\text{ m/s}.
        \]
        Therefore, the distance between the two particles is increasing at a rate of $5\text{ m/s}$.
\end{soaljawab}

\begin{soaljawab}{3}
\item Draw the graph of $y = |2x - 2|$ for $0 \le x \le 3$ and calculate
        \[
          \int_{0}^{3} |2x - 2| \, dx
        \]
        using the appropriate formula from geometry.
\jawaban
The function $y = |2x - 2| = 2|x - 1|$ forms two right triangles on $[0, 3]$ with vertex on the $x$-axis at $x = 1$:
        \begin{enumerate}
          \item Triangle 1 on $[0, 1]$:
                Base $b_1 = 1 - 0 = 1$, height $h_1 = |2(0) - 2| = 2$.
                $\text{Area}_1 = \frac{1}{2} \cdot 1 \cdot 2 = 1$.
          \item Triangle 2 on $[1, 3]$:
                Base $b_2 = 3 - 1 = 2$, height $h_2 = |2(3) - 2| = 4$.
                $\text{Area}_2 = \frac{1}{2} \cdot 2 \cdot 4 = 4$.
        \end{enumerate}
        Using plane geometry:
        \[
          \int_{0}^{3} |2x - 2| \, dx = \text{Area}_1 + \text{Area}_2 = 1 + 4 = 5.
        \]
\end{soaljawab}

\begin{soaljawab}{4}
\item Given $z_1 = 5 - 10i$, $z_2 = -3 + 2i$, and $z_3 = 3 + 4i$.
        \begin{enumerate}
          \item Calculate $\dfrac{z_1 z_2}{z_3}$ and express it in the form $a + bi$.
          \item Calculate $\left| \dfrac{z_1 z_2}{z_3} \right|^2$.
        \end{enumerate}
\jawaban
\begin{enumerate}
          \item Compute $z_1 z_2$:
                \[
                  z_1 z_2 = (5 - 10i)(-3 + 2i) = -15 + 10i + 30i - 20i^2 = -15 + 40i + 20 = 5 + 40i.
                \]
                Now divide by $z_3 = 3 + 4i$:
                \begin{align*}
                  \frac{z_1 z_2}{z_3} & = \frac{5 + 40i}{3 + 4i} \cdot \frac{3 - 4i}{3 - 4i} = \frac{15 - 20i + 120i - 160i^2}{3^2 + 4^2} \\
                                      & = \frac{15 + 100i + 160}{25} = \frac{175 + 100i}{25} = 7 + 4i.
                \end{align*}
          \item Modulus squared:
                \[
                  \left| \frac{z_1 z_2}{z_3} \right|^2 = |7 + 4i|^2 = 7^2 + 4^2 = 49 + 16 = 65.
                \]
        \end{enumerate}
\end{soaljawab}

\begin{soaljawab}{5}
\item Given the following system of linear equations
        \[
          \begin{aligned}
            x + y + 2z   & = 8,  \\
            -x - 2y + 3z & = 1,  \\
            3x - 7y + 4z & = 10.
          \end{aligned}
        \]
        \begin{enumerate}
          \item Convert the system into an augmented matrix.
          \item Solve it using the Gauss elimination method or the Gauss--Jordan elimination method.
        \end{enumerate}
\jawaban
\begin{enumerate}
\item The columns correspond to $x$, $y$, $z$, and the right-hand side, respectively.
The augmented matrix is
\[
    [A\mid\mathbf b] = \left[\begin{array}{ccc|c}
            1 & 1 & 2 & 8 \\
            -1 & -2 & 3 & 1 \\
            3 & -7 & 4 & 10
        \end{array}\right].
\]
\item Apply Gaussian elimination to obtain an upper triangular coefficient matrix.
Here $B_i$ denotes row $i$; an arrow labelled $B_i+kB_j$ replaces row $i$
by $B_i+kB_j$.
{\setlength{\arraycolsep}{4bp}
\begin{align*}
    & \left[\begin{array}{ccc|c}
            1 & 1 & 2 & 8 \\
            -1 & -2 & 3 & 1 \\
            3 & -7 & 4 & 10
        \end{array}\right]
        \begin{array}{c}\xrightarrow{B_2+B_1}\end{array}
        \left[\begin{array}{ccc|c}
            1 & 1 & 2 & 8 \\
            0 & -1 & 5 & 9 \\
            3 & -7 & 4 & 10
        \end{array}\right] \\[6bp]
    & \begin{array}{c}\xrightarrow{B_3-3B_1}\end{array}
        \left[\begin{array}{ccc|c}
            1 & 1 & 2 & 8 \\
            0 & -1 & 5 & 9 \\
            0 & -10 & -2 & -14
        \end{array}\right]
        \begin{array}{c}\xrightarrow{B_3-10B_2}\end{array}
        \left[\begin{array}{ccc|c}
            1 & 1 & 2 & 8 \\
            0 & -1 & 5 & 9 \\
            0 & 0 & -52 & -104
        \end{array}\right]
\end{align*}}
Now use back substitution, starting with the last row:
\begin{align*}
    B_3:\quad -52z&=-104
        &\Longrightarrow\quad z&=2,\\
    B_2:\quad -y+5(2)&=9
        &\Longrightarrow\quad -y&=-1
        \quad\Longrightarrow\quad y=1,\\
    B_1:\quad x+1+2(2)&=8
        &\Longrightarrow\quad x&=8-1-4=3.
\end{align*}
Therefore, $(x,y,z)=(3,1,2)$.
\end{enumerate}
\end{soaljawab}

\label{akhir-paket-2}

% PAKET 3: Senin, 8 Desember 2025 | 11.00 - 12.40 WIB | Reguler (Kelas 17-22)
\startpacket{3}
\pdfbookmark[0]{Paket 3: Senin 11.00 Reguler (Kelas 17--22)}{paket3}
\examheaderID{Kalkulus 1, 3 SKS}{Senin, 8 Desember 2025}{11:00 -- 12:40, 100 menit}{17 -- 22, Tertutup}
\warningsID
\cpmktableID
\scoretableID
\begin{soaljawab}{1}
\item Hitunglah integral berikut
        \[
          \int \frac{\sin(7x)}{(1 + \cos(7x))^3} \, dx.
        \]
\jawaban
Gunakan substitusi $u = 1 + \cos(7x)$, sehingga $du = -7\sin(7x)\,dx \implies \sin(7x)\,dx = -\frac{1}{7}\,du$.
        \begin{align*}
          \int \frac{\sin(7x)}{(1 + \cos(7x))^3} \, dx & = \int u^{-3} \left(-\frac{1}{7}\right) du = -\frac{1}{7} \left(\frac{u^{-2}}{-2}\right) + C \\
                                                       & = \frac{1}{14(1 + \cos(7x))^2} + C.
        \end{align*}
\end{soaljawab}

\begin{soaljawab}{2}
\item Diberikan $y = 2x^3 - 3x^2$.
        \begin{enumerate}
          \item Carilah perpotongan fungsi tersebut dengan sumbu $y$ dan sumbu $x$.
          \item Carilah turunan pertama, titik stasioner, dan selang naik turun dari fungsi tersebut.
          \item Carilah turunan kedua, titik belok, dan kecekungan dari fungsi tersebut.
          \item Gunakan informasi di atas untuk mendapatkan grafik dari fungsi tersebut.
        \end{enumerate}
\jawaban
\begin{enumerate}
          \item Titik potong sumbu:
                \begin{itemize}
                  \item Sumbu $y$ ($x = 0$): $y = 0 \implies (0, 0)$.
                  \item Sumbu $x$ ($y = 0$): $2x^3 - 3x^2 = 0 \implies x^2(2x - 3) = 0 \implies x = 0$ atau $x = \frac{3}{2}$. Jadi titik potongnya $(0, 0)$ dan $\left(\frac{3}{2}, 0\right)$.
                \end{itemize}
          \item Turunan pertama:
                \[
                  y' = 6x^2 - 6x = 6x(x - 1).
                \]
                Titik stasioner saat $y' = 0 \implies x = 0$ dan $x = 1$.
                \begin{itemize}
                  \item Naik pada $(-\infty, 0) \cup (1, \infty)$.
                  \item Turun pada $(0, 1)$.
                \end{itemize}
          \item Turunan kedua:
                \[
                  y'' = 12x - 6 = 6(2x - 1).
                \]
                Titik belok saat $y'' = 0 \implies x = \frac{1}{2}$.
                Nilai $y\left(\frac{1}{2}\right) = 2\left(\frac{1}{8}\right) - 3\left(\frac{1}{4}\right) = \frac{1}{4} - \frac{3}{4} = -\frac{1}{2}$. Titik belok di $\left(\frac{1}{2}, -\frac{1}{2}\right)$.
                \begin{itemize}
                  \item Cekung ke bawah pada $\left(-\infty, \frac{1}{2}\right)$.
                  \item Cekung ke atas pada $\left(\frac{1}{2}, \infty\right)$.
                \end{itemize}
          \item Sketsa kurva memiliki titik maksimum lokal di $(0, 0)$, titik minimum lokal di $(1, -1)$, dan titik belok di $\left(\frac{1}{2}, -\frac{1}{2}\right)$.
        \end{enumerate}
\end{soaljawab}

\begin{soaljawab}{3}
\item Diketahui $\displaystyle \int_{1}^{8} f(x) \, dx = 4$. Dapatkan
        \[
          \int_{1}^{2} x^2 f(9 - x^3) \, dx
        \]
        dengan substitusi $u = 9 - x^3$.
\jawaban
Substitusi $u = 9 - x^3$, maka $du = -3x^2\,dx \implies x^2\,dx = -\frac{1}{3}\,du$.
        Batas integrasi:
        \begin{itemize}
          \item Saat $x = 1 \implies u = 9 - 1^3 = 8$.
          \item Saat $x = 2 \implies u = 9 - 2^3 = 1$.
        \end{itemize}
        Maka:
        \begin{align*}
          \int_{1}^{2} x^2 f(9 - x^3) \, dx & = \int_{8}^{1} f(u) \left(-\frac{1}{3}\right) du = \frac{1}{3} \int_{1}^{8} f(u) \, du \\
                                            & = \frac{1}{3} (4) = \frac{4}{3}.
        \end{align*}
\end{soaljawab}

\begin{soaljawab}{4}
\item Misalkan $z = 3 + 3i$.
        \begin{enumerate}
          \item Nyatakan $z$ dalam bentuk kutub.
          \item Hitunglah $z^5$ dalam bentuk kutub.
          \item Nyatakan $z^5$ dalam bentuk $a + bi$.
        \end{enumerate}
\jawaban
\begin{enumerate}
          \item Modulus $r = \sqrt{3^2 + 3^2} = \sqrt{18} = 3\sqrt{2}$.
                Argumen $\theta = \arctan\left(\frac{3}{3}\right) = 45^\circ = \frac{\pi}{4}$.
                Bentuk kutub: $z = 3\sqrt{2}\operatorname{cis}(45^\circ)$.
          \item Menggunakan Teorema De Moivre:
                \begin{align*}
                  z^5 & = (3\sqrt{2})^5 \operatorname{cis}(5 \times 45^\circ) = 3^5 \cdot 2^{5/2} \operatorname{cis}(225^\circ) \\
                      & = 243 \cdot 4\sqrt{2} \operatorname{cis}(225^\circ) = 972\sqrt{2} \operatorname{cis}(225^\circ).
                \end{align*}
          \item Bentuk $a + bi$:
                \begin{align*}
                  z^5 & = 972\sqrt{2} \left(\cos 225^\circ + i\sin 225^\circ\right) = 972\sqrt{2} \left(-\frac{\sqrt{2}}{2} - i\frac{\sqrt{2}}{2}\right) \\
                      & = 972\sqrt{2}\left(-\frac{\sqrt{2}}{2}\right) - i 972\sqrt{2}\left(\frac{\sqrt{2}}{2}\right) = -972 - 972i.
                \end{align*}
        \end{enumerate}
\end{soaljawab}

\begin{soaljawab}{5}
\item Carilah invers dari
        \[
          A = \begin{bmatrix}
            1 & 2 & 3  \\
            1 & 3 & 6  \\
            2 & 6 & 13
          \end{bmatrix}
        \]
        dengan menggunakan operasi baris elementer.
\jawaban
Gabungkan $A$ dengan matriks identitas $I_3$. Lakukan OBE pada
kedua bagian hingga $[A\mid I_3]$ menjadi $[I_3\mid A^{-1}]$.
Di sini $B_i$ menyatakan baris ke-$i$; label $B_i+kB_j$ berarti
mengganti $B_i$ dengan $B_i+kB_j$.
{\setlength{\arraycolsep}{4bp}
\begin{align*}
    & \left[\begin{array}{ccc|ccc}
            1 & 2 & 3 & 1 & 0 & 0 \\
            1 & 3 & 6 & 0 & 1 & 0 \\
            2 & 6 & 13 & 0 & 0 & 1
        \end{array}\right]
        \begin{array}{c}\xrightarrow{B_2-B_1}\end{array}
        \left[\begin{array}{ccc|ccc}
            1 & 2 & 3 & 1 & 0 & 0 \\
            0 & 1 & 3 & -1 & 1 & 0 \\
            2 & 6 & 13 & 0 & 0 & 1
        \end{array}\right] \\[6bp]
    & \begin{array}{c}\xrightarrow{B_3-2B_1}\end{array}
        \left[\begin{array}{ccc|ccc}
            1 & 2 & 3 & 1 & 0 & 0 \\
            0 & 1 & 3 & -1 & 1 & 0 \\
            0 & 2 & 7 & -2 & 0 & 1
        \end{array}\right]
        \begin{array}{c}\xrightarrow{B_3-2B_2}\end{array}
        \left[\begin{array}{ccc|ccc}
            1 & 2 & 3 & 1 & 0 & 0 \\
            0 & 1 & 3 & -1 & 1 & 0 \\
            0 & 0 & 1 & 0 & -2 & 1
        \end{array}\right] \\[6bp]
    & \begin{array}{c}\xrightarrow{B_1-2B_2}\end{array}
        \left[\begin{array}{ccc|ccc}
            1 & 0 & -3 & 3 & -2 & 0 \\
            0 & 1 & 3 & -1 & 1 & 0 \\
            0 & 0 & 1 & 0 & -2 & 1
        \end{array}\right]
        \begin{array}{c}\xrightarrow{B_1+3B_3}\end{array}
        \left[\begin{array}{ccc|ccc}
            1 & 0 & 0 & 3 & -8 & 3 \\
            0 & 1 & 3 & -1 & 1 & 0 \\
            0 & 0 & 1 & 0 & -2 & 1
        \end{array}\right] \\[6bp]
    & \begin{array}{c}\xrightarrow{B_2-3B_3}\end{array}
        \left[\begin{array}{ccc|ccc}
            1 & 0 & 0 & 3 & -8 & 3 \\
            0 & 1 & 0 & -1 & 7 & -3 \\
            0 & 0 & 1 & 0 & -2 & 1
        \end{array}\right]
\end{align*}}
Bagian kiri sudah menjadi $I_3$, sehingga bagian kanan merupakan inversnya:
\[
    A^{-1} = \begin{bmatrix}
            3 & -8 & 3 \\
            -1 & 7 & -3 \\
            0 & -2 & 1
        \end{bmatrix}.
\]
\end{soaljawab}

\label{akhir-paket-3}

% PAKET 4: Senin, 8 Desember 2025 | 13.30 - 15.10 WIB | Reguler (Kelas 13-16, 50-61)
\startpacket{4}
\pdfbookmark[0]{Paket 4: Senin 13.30 Reguler (Kelas 13--16, 50--61)}{paket4}
\examheaderID{Kalkulus 1, 3 SKS}{Senin, 8 Desember 2025}{13:30 -- 15:10, 100 menit}{13 -- 16 dan 50 -- 61, Tertutup}
\warningsID
\cpmktableID
\scoretableID
\begin{soaljawab}{1}
\item Hitunglah integral berikut
        \[
          \int \frac{x^2}{(x^3 + 2)^{1/4}} \, dx.
        \]
\jawaban
Substitusi $u = x^3 + 2 \implies du = 3x^2\,dx \implies x^2\,dx = \frac{1}{3}\,du$.
        \begin{align*}
          \int \frac{x^2}{(x^3 + 2)^{1/4}} \, dx & = \frac{1}{3} \int u^{-1/4} \, du = \frac{1}{3} \cdot \frac{4}{3} u^{3/4} + C \\
                                                 & = \frac{4}{9} (x^3 + 2)^{3/4} + C.
        \end{align*}
\end{soaljawab}

\begin{soaljawab}{2}
\item Diberikan fungsi $y = 2x + \dfrac{1}{(x-1)^2}$ dengan $x \in (1, \infty)$.
        \begin{enumerate}
          \item Tentukan maksimum atau minimum relatif (jika ada) dari $y$.
          \item Apakah hasil pada bagian (a) juga merupakan maksimum atau minimum absolut?
        \end{enumerate}
\jawaban
\begin{enumerate}
          \item Cari turunan pertama untuk $x > 1$:
                \[
                  y' = 2 - 2(x-1)^{-3} = 2 - \frac{2}{(x-1)^3}.
                \]
                Titik kritis: $y' = 0 \implies \frac{2}{(x-1)^3} = 2 \implies (x-1)^3 = 1 \implies x - 1 = 1 \implies x = 2$.
                Turunan kedua:
                \[
                  y'' = 6(x-1)^{-4} = \frac{6}{(x-1)^4} > 0 \quad \forall x \in (1, \infty).
                \]
                Karena $y''(2) = 6 > 0$, maka $x = 2$ memberikan minimum relatif dengan nilai:
                \[
                  y(2) = 2(2) + \frac{1}{(2-1)^2} = 4 + 1 = 5.
                \]
          \item Karena $y'' > 0$ untuk seluruh domain $(1, \infty)$ (fungsi konveks murni), dan $\displaystyle \lim_{x \to 1^+} y = \infty$ serta $\displaystyle \lim_{x \to \infty} y = \infty$, maka nilai $y(2) = 5$ adalah minimum absolut pada interval $(1, \infty)$.
        \end{enumerate}
\end{soaljawab}

\begin{soaljawab}{3}
\item Diketahui $\displaystyle \int_{0}^{4} f(x) \, dx = 3$. Dapatkan
        \[
          \int_{0}^{16} \frac{f(\sqrt{x})}{\sqrt{x}} \, dx
        \]
        dengan substitusi $u = \sqrt{x}$.
\jawaban
Substitusi $u = \sqrt{x} \implies du = \frac{1}{2\sqrt{x}}\,dx \implies \frac{dx}{\sqrt{x}} = 2\,du$.
        Batas integrasi:
        \begin{itemize}
          \item $x = 0 \implies u = 0$
          \item $x = 16 \implies u = \sqrt{16} = 4$
        \end{itemize}
        Maka:
        \[
          \int_{0}^{16} \frac{f(\sqrt{x})}{\sqrt{x}} \, dx = \int_{0}^{4} f(u) (2\,du) = 2 \int_{0}^{4} f(u)\,du = 2(3) = 6.
        \]
\end{soaljawab}

\begin{soaljawab}{4}
\item Nyatakan bilangan kompleks $2 + 2\sqrt{3}i$ dalam bentuk kutub. Carilah semua bilangan kompleks $z$ yang memenuhi
        \[
          z^3 = 2 + 2\sqrt{3}i.
        \]
\jawaban
Misalkan $w = 2 + 2\sqrt{3}i$.
        Modulus $r = \sqrt{2^2 + (2\sqrt{3})^2} = \sqrt{4 + 12} = \sqrt{16} = 4$.
        Argumen $\theta = \arctan\left(\frac{2\sqrt{3}}{2}\right) = \arctan(\sqrt{3}) = 60^\circ = \frac{\pi}{3}$.
        Bentuk kutub: $w = 4\operatorname{cis}(60^\circ)$.

        Penyelesaian $z^3 = w$ untuk $k = 0, 1, 2$:
        \[
          z_k = \sqrt[3]{4} \operatorname{cis}\left(\frac{60^\circ + 360^\circ k}{3}\right) = 4^{1/3} \operatorname{cis}(20^\circ + 120^\circ k).
        \]
        Diperoleh:
        \begin{align*}
          z_0 & = 4^{1/3} \operatorname{cis}(20^\circ),  \\
          z_1 & = 4^{1/3} \operatorname{cis}(140^\circ), \\
          z_2 & = 4^{1/3} \operatorname{cis}(260^\circ).
        \end{align*}
\end{soaljawab}

\begin{soaljawab}{5}
\item Carilah invers dari
        \[
          A = \begin{bmatrix}
            -1 & -2 & -2 \\
            1  & 3  & 3  \\
            -1 & 4  & 5
          \end{bmatrix}
        \]
        dengan menggunakan operasi baris elementer.
\jawaban
Gabungkan $A$ dengan matriks identitas $I_3$. Lakukan OBE pada
kedua bagian hingga $[A\mid I_3]$ menjadi $[I_3\mid A^{-1}]$.
Di sini $B_i$ menyatakan baris ke-$i$; label $B_i+kB_j$ berarti
mengganti $B_i$ dengan $B_i+kB_j$.
{\setlength{\arraycolsep}{4bp}
\begin{align*}
    & \left[\begin{array}{ccc|ccc}
            -1 & -2 & -2 & 1 & 0 & 0 \\
            1 & 3 & 3 & 0 & 1 & 0 \\
            -1 & 4 & 5 & 0 & 0 & 1
        \end{array}\right]
        \begin{array}{c}\xrightarrow{-B_1}\end{array}
        \left[\begin{array}{ccc|ccc}
            1 & 2 & 2 & -1 & 0 & 0 \\
            1 & 3 & 3 & 0 & 1 & 0 \\
            -1 & 4 & 5 & 0 & 0 & 1
        \end{array}\right] \\[6bp]
    & \begin{array}{c}\xrightarrow{B_2-B_1}\end{array}
        \left[\begin{array}{ccc|ccc}
            1 & 2 & 2 & -1 & 0 & 0 \\
            0 & 1 & 1 & 1 & 1 & 0 \\
            -1 & 4 & 5 & 0 & 0 & 1
        \end{array}\right]
        \begin{array}{c}\xrightarrow{B_3+B_1}\end{array}
        \left[\begin{array}{ccc|ccc}
            1 & 2 & 2 & -1 & 0 & 0 \\
            0 & 1 & 1 & 1 & 1 & 0 \\
            0 & 6 & 7 & -1 & 0 & 1
        \end{array}\right] \\[6bp]
    & \begin{array}{c}\xrightarrow{B_3-6B_2}\end{array}
        \left[\begin{array}{ccc|ccc}
            1 & 2 & 2 & -1 & 0 & 0 \\
            0 & 1 & 1 & 1 & 1 & 0 \\
            0 & 0 & 1 & -7 & -6 & 1
        \end{array}\right]
        \begin{array}{c}\xrightarrow{B_1-2B_2}\end{array}
        \left[\begin{array}{ccc|ccc}
            1 & 0 & 0 & -3 & -2 & 0 \\
            0 & 1 & 1 & 1 & 1 & 0 \\
            0 & 0 & 1 & -7 & -6 & 1
        \end{array}\right] \\[6bp]
    & \begin{array}{c}\xrightarrow{B_2-B_3}\end{array}
        \left[\begin{array}{ccc|ccc}
            1 & 0 & 0 & -3 & -2 & 0 \\
            0 & 1 & 0 & 8 & 7 & -1 \\
            0 & 0 & 1 & -7 & -6 & 1
        \end{array}\right]
\end{align*}}
Bagian kiri sudah menjadi $I_3$, sehingga bagian kanan merupakan inversnya:
\[
    A^{-1} = \begin{bmatrix}
            -3 & -2 & 0 \\
            8 & 7 & -1 \\
            -7 & -6 & 1
        \end{bmatrix}.
\]
\end{soaljawab}

\label{akhir-paket-4}

% PAKET 5: Selasa, 9 Desember 2025 | 07.00 - 08.40 WIB | IUP (Class 2, 3, 68, 69)
\startpacket{5}
\pdfbookmark[0]{Paket 5: Tuesday 07.00 IUP (Class 2, 3, 68, 69)}{paket5}
\examheaderEN{Kalkulus 1, 3 SKS}{Tuesday, 9 December 2025}{07:00 -- 08:40, 100 minutes}{2, 3, 68, 69, Closed book}
\warningsEN
\cpmktableEN
\scoretableEN
\begin{soaljawab}{1}
\item Calculate the following integral
        \[
          \int \frac{x}{(4x^2 + 1)^3} \, dx.
        \]
\jawaban
Let $u = 4x^2 + 1 \implies du = 8x\,dx \implies x\,dx = \frac{1}{8}\,du$.
        \begin{align*}
          \int \frac{x}{(4x^2 + 1)^3} \, dx & = \frac{1}{8} \int u^{-3} \, du = \frac{1}{8} \left(\frac{u^{-2}}{-2}\right) + C \\
                                            & = -\frac{1}{16(4x^2 + 1)^2} + C.
        \end{align*}
\end{soaljawab}

\begin{soaljawab}{2}
\item A tank in the shape of an inverted cone has a height of $12\text{ m}$ and a circular radius of $6\text{ m}$. Water is poured into the tank at a rate of $2\text{ m}^3/\text{min}$. Determine the rate of change of the height of water level when the height of the water (measured from the bottom) reaches $6\text{ m}$. Use the volume formula for a cone $V = \frac{1}{3}\pi r^2 y$.
\jawaban
For an inverted cone with height $H = 12$ and radius $R = 6$, by similar triangles:
        \[
          \frac{r}{y} = \frac{6}{12} = \frac{1}{2} \implies r = \frac{y}{2}.
        \]
        Volume in terms of height $y$:
        \[
          V = \frac{1}{3}\pi r^2 y = \frac{1}{3}\pi \left(\frac{y}{2}\right)^2 y = \frac{\pi y^3}{12}.
        \]
        Differentiating with respect to time $t$:
        \[
          \frac{dV}{dt} = \frac{\pi y^2}{4} \frac{dy}{dt} \implies \frac{dy}{dt} = \frac{4}{\pi y^2} \frac{dV}{dt}.
        \]
        Given $\dfrac{dV}{dt} = 2\text{ m}^3/\text{min}$ and $y = 6\text{ m}$:
        \[
          \frac{dy}{dt} = \frac{4}{\pi (6^2)} \cdot 2 = \frac{8}{36\pi} = \frac{2}{9\pi}\text{ m/min}.
        \]
\end{soaljawab}

\begin{soaljawab}{3}
\item Draw the graph of $y = 1 - \dfrac{x}{2}$ for $0 \le x \le 6$ and calculate
        \[
          \int_{0}^{6} \left(1 - \frac{x}{2}\right) \, dx
        \]
        using the appropriate formula from geometry.
\jawaban
The line $y = 1 - \frac{x}{2}$ crosses the $x$-axis at $x = 2$:
        \begin{itemize}
          \item On $[0, 2]$, $y \ge 0$. It forms a triangle above the $x$-axis with base $2$ and height $1$:
                \[
                  A_1 = \frac{1}{2} \cdot 2 \cdot 1 = 1.
                \]
          \item On $[2, 6]$, $y \le 0$. It forms a triangle below the $x$-axis with base $6 - 2 = 4$ and height $|1 - 3| = 2$:
                \[
                  A_2 = \frac{1}{2} \cdot 4 \cdot 2 = 4.
                \]
        \end{itemize}
        Using net signed area:
        \[
          \int_{0}^{6} \left(1 - \frac{x}{2}\right) \, dx = A_1 - A_2 = 1 - 4 = -3.
        \]
\end{soaljawab}

\begin{soaljawab}{4}
\item Given $w = \dfrac{10 + 2i}{3 - 2i}$.
        \begin{enumerate}
          \item Express $w$ in the form $w = a + bi$.
          \item Express $w$ in polar form.
        \end{enumerate}
\jawaban
\begin{enumerate}
          \item Simplify $w$:
                \begin{align*}
                  w & = \frac{10 + 2i}{3 - 2i} \cdot \frac{3 + 2i}{3 + 2i} = \frac{30 + 20i + 6i + 4i^2}{3^2 + (-2)^2} \\
                    & = \frac{26 + 26i}{13} = 2 + 2i.
                \end{align*}
          \item Polar form:
                $r = \sqrt{2^2 + 2^2} = \sqrt{8} = 2\sqrt{2}$.
                $\theta = \arctan(2/2) = 45^\circ = \frac{\pi}{4}$.
                \[
                  w = 2\sqrt{2} \operatorname{cis}(45^\circ) = 2\sqrt{2}\left(\cos\frac{\pi}{4} + i\sin\frac{\pi}{4}\right).
                \]
        \end{enumerate}
\end{soaljawab}

\begin{soaljawab}{5}
\item Given the following system of linear equations
        \[
          \begin{aligned}
            x + y + 2z    & = 4,  \\
            2x + 3y + 6z  & = 10, \\
            3x + 6y + 10z & = 17.
          \end{aligned}
        \]
        \begin{enumerate}
          \item Convert the system into an augmented matrix.
          \item Solve it using the Gauss elimination method or the Gauss--Jordan elimination method.
        \end{enumerate}
\jawaban
\begin{enumerate}
\item The columns correspond to $x$, $y$, $z$, and the right-hand side, respectively.
The augmented matrix is
\[
    [A\mid\mathbf b] = \left[\begin{array}{ccc|c}
            1 & 1 & 2 & 4 \\
            2 & 3 & 6 & 10 \\
            3 & 6 & 10 & 17
        \end{array}\right].
\]
\item Apply Gaussian elimination to obtain an upper triangular coefficient matrix.
Here $B_i$ denotes row $i$; an arrow labelled $B_i+kB_j$ replaces row $i$
by $B_i+kB_j$.
{\setlength{\arraycolsep}{4bp}
\begin{align*}
    & \left[\begin{array}{ccc|c}
            1 & 1 & 2 & 4 \\
            2 & 3 & 6 & 10 \\
            3 & 6 & 10 & 17
        \end{array}\right]
        \begin{array}{c}\xrightarrow{B_2-2B_1}\end{array}
        \left[\begin{array}{ccc|c}
            1 & 1 & 2 & 4 \\
            0 & 1 & 2 & 2 \\
            3 & 6 & 10 & 17
        \end{array}\right] \\[6bp]
    & \begin{array}{c}\xrightarrow{B_3-3B_1}\end{array}
        \left[\begin{array}{ccc|c}
            1 & 1 & 2 & 4 \\
            0 & 1 & 2 & 2 \\
            0 & 3 & 4 & 5
        \end{array}\right]
        \begin{array}{c}\xrightarrow{B_3-3B_2}\end{array}
        \left[\begin{array}{ccc|c}
            1 & 1 & 2 & 4 \\
            0 & 1 & 2 & 2 \\
            0 & 0 & -2 & -1
        \end{array}\right]
\end{align*}}
Use back substitution from the third row to the first:
\begin{align*}
    B_3:\quad -2z&=-1
        &\Longrightarrow\quad z&=\frac12,\\
    B_2:\quad y+2\left(\frac12\right)&=2
        &\Longrightarrow\quad y&=2-1=1,\\
    B_1:\quad x+1+2\left(\frac12\right)&=4
        &\Longrightarrow\quad x&=4-1-1=2.
\end{align*}
Therefore, $(x,y,z)=\left(2,1,\frac12\right)$.
\end{enumerate}
\end{soaljawab}

\label{akhir-paket-5}

% PAKET 6: Selasa, 9 Desember 2025 | 07.00 - 08.40 WIB | Mengulang (Kelas 113, 114)
\startpacket{6}
\pdfbookmark[0]{Paket 6: Selasa 07.00 Mengulang (Kelas 113, 114)}{paket6}
\examheaderID{Kalkulus 1, 3 SKS}{Selasa, 9 Desember 2025}{07:00 -- 08:40, 100 menit}{113, 114, Tertutup}
\warningsID
\cpmktableID
\scoretableID
\begin{soaljawab}{1}
\item Hitunglah integral berikut
        \[
          \int 5x^4(x^5 + 3)^6 \, dx.
        \]
\jawaban
Substitusi $u = x^5 + 3 \implies du = 5x^4\,dx$:
        \[
          \int 5x^4(x^5 + 3)^6 \, dx = \int u^6 \, du = \frac{1}{7} u^7 + C = \frac{1}{7} (x^5 + 3)^7 + C.
        \]
\end{soaljawab}

\begin{soaljawab}{2}
\item Dapatkan turunan pertama, titik stasioner, dan maksimum dan minimum relatif dari
        \[
          f(x) = x^4 - 2x^2 - 15.
        \]
\jawaban
\begin{itemize}
          \item Turunan pertama:
                \[
                  f'(x) = 4x^3 - 4x = 4x(x^2 - 1) = 4x(x - 1)(x + 1).
                \]
          \item Titik stasioner: $f'(x) = 0 \implies x = -1, 0, 1$.
          \item Uji turunan kedua: $f''(x) = 12x^2 - 4$.
                \begin{itemize}
                  \item $x = 0$: $f''(0) = -4 < 0 \implies$ maksimum relatif di $(0, f(0)) = (0, -15)$.
                  \item $x = -1$: $f''(-1) = 8 > 0 \implies$ minimum relatif di $(-1, f(-1)) = (-1, -16)$.
                  \item $x = 1$: $f''(1) = 8 > 0 \implies$ minimum relatif di $(1, f(1)) = (1, -16)$.
                \end{itemize}
        \end{itemize}
\end{soaljawab}

\begin{soaljawab}{3}
\item Gambarkan grafik dari $y = 2x + 1$ untuk $1 \le x \le 3$ dan hitung integral
        \[
          \int_{1}^{3} (2x + 1) \, dx
        \]
        dengan menggunakan rumus dari geometri bidang yang sesuai.
\jawaban
Daerah di bawah garis $y = 2x + 1$ pada interval $[1, 3]$ merupakan trapesium tegak dengan:
        \begin{itemize}
          \item Sisi sejajar $y_1 = 2(1) + 1 = 3$ dan $y_2 = 2(3) + 1 = 7$
          \item Tinggi trapesium $t = 3 - 1 = 2$
        \end{itemize}
        Luas trapesium:
        \[
          \int_{1}^{3} (2x + 1) \, dx = \frac{y_1 + y_2}{2} \cdot t = \frac{3 + 7}{2} \cdot 2 = 10.
        \]
\end{soaljawab}

\begin{soaljawab}{4}
\item Diberikan bilangan kompleks $z_1 = 1 + 2i$, $z_2 = 2 + 3i$, dan $z_3 = 3 + 4i$. Sederhanakan
        \[
          \frac{5(z_1 + z_2 + z_3)}{z_1 \times z_2}
        \]
        dalam bentuk $a + bi$.
\jawaban
Pembilang:
        \[
          z_1 + z_2 + z_3 = (1 + 2i) + (2 + 3i) + (3 + 4i) = 6 + 9i.
        \]
        \[
          5(z_1 + z_2 + z_3) = 5(6 + 9i) = 30 + 45i.
        \]
        Penyebut:
        \[
          z_1 z_2 = (1 + 2i)(2 + 3i) = 2 + 3i + 4i + 6i^2 = -4 + 7i.
        \]
        Pembagian:
        \begin{align*}
          \frac{30 + 45i}{-4 + 7i} & = \frac{30 + 45i}{-4 + 7i} \cdot \frac{-4 - 7i}{-4 - 7i} = \frac{-120 - 210i - 180i - 315i^2}{(-4)^2 + 7^2} \\
                                   & = \frac{-120 - 390i + 315}{16 + 49} = \frac{195 - 390i}{65} = 3 - 6i.
        \end{align*}
\end{soaljawab}

\begin{soaljawab}{5}
\item Misalkan $A = \begin{bmatrix} x & y \\ 0 & z \end{bmatrix}$ dan $A^3 = \begin{bmatrix} 8 & -57 \\ 0 & 27 \end{bmatrix}$. Carilah $x, y$, dan $z$.
\jawaban
\textbf{Langkah 1: hitung $A^2$ dengan perkalian baris dan kolom.}
\begin{align*}
    A^2
    &=\begin{bmatrix}x&y\\0&z\end{bmatrix}
      \begin{bmatrix}x&y\\0&z\end{bmatrix}\\
    &=\begin{bmatrix}
        x\cdot x+y\cdot0 & x\cdot y+y\cdot z\\
        0\cdot x+z\cdot0 & 0\cdot y+z\cdot z
      \end{bmatrix}
     =\begin{bmatrix}x^2&y(x+z)\\0&z^2\end{bmatrix}.
\end{align*}
\textbf{Langkah 2: kalikan $A^2$ dengan $A$.}
\begin{align*}
    A^3
    &=\begin{bmatrix}x^2&y(x+z)\\0&z^2\end{bmatrix}
      \begin{bmatrix}x&y\\0&z\end{bmatrix}\\
    &=\begin{bmatrix}
        x^2\cdot x+y(x+z)\cdot0 & x^2\cdot y+y(x+z)\cdot z\\
        0\cdot x+z^2\cdot0 & 0\cdot y+z^2\cdot z
      \end{bmatrix}\\
    &=\begin{bmatrix}x^3&x^2y+xyz+yz^2\\0&z^3\end{bmatrix}
     =\begin{bmatrix}x^3&y(x^2+xz+z^2)\\0&z^3\end{bmatrix}.
\end{align*}
\textbf{Langkah 3: samakan elemen yang bersesuaian.}
Karena $A^3=\begin{bmatrix}8&-57\\0&27\end{bmatrix}$, diperoleh
\[
    x^3=8,\qquad z^3=27,\qquad y(x^2+xz+z^2)=-57.
\]
Untuk $x,y,z$ real, dua persamaan pertama memberi $x=2$ dan $z=3$.
Substitusikan ke persamaan ketiga:
\begin{align*}
    y(2^2+2\cdot3+3^2)&=-57,\\
    y(4+6+9)&=-57,\\
    19y&=-57
    \quad\Longrightarrow\quad y=-3.
\end{align*}
Jadi $x=2$, $y=-3$, dan $z=3$.
\end{soaljawab}

\label{akhir-paket-6}

% PAKET 7: Selasa, 9 Desember 2025 | 09.00 - 10.40 WIB | Reguler (Kelas 5-12, 23-29)
\startpacket{7}
\pdfbookmark[0]{Paket 7: Selasa 09.00 Reguler (Kelas 5--12, 23--29)}{paket7}
\examheaderID{Kalkulus 1, 3 SKS}{Selasa, 9 Desember 2025}{09:00 -- 10:40, 100 menit}{5 -- 12, 23 -- 29, Tertutup}
\warningsID
\cpmktableID
\scoretableID
\begin{soaljawab}{1}
\item Hitunglah integral berikut
        \[
          \int \frac{3x^3}{(3x^4 + 2)^{1/3}} \, dx.
        \]
\jawaban
Substitusi $u = 3x^4 + 2 \implies du = 12x^3\,dx \implies 3x^3\,dx = \frac{1}{4}\,du$.
        \begin{align*}
          \int \frac{3x^3}{(3x^4 + 2)^{1/3}} \, dx & = \frac{1}{4} \int u^{-1/3} \, du = \frac{1}{4} \cdot \frac{3}{2} u^{2/3} + C \\
                                                   & = \frac{3}{8} (3x^4 + 2)^{2/3} + C.
        \end{align*}
\end{soaljawab}

\begin{soaljawab}{2}
\item Diberikan $y = -2x^4 + 8x^2$.
        \begin{enumerate}
          \item Carilah turunan pertama, titik stasioner, maksimum dan minimum relatif, dan selang naik turun dari fungsi tersebut.
          \item Carilah turunan kedua, titik belok, dan kecekungan dari fungsi tersebut.
        \end{enumerate}
\jawaban
\begin{enumerate}
          \item Turunan pertama:
                \[
                  y' = -8x^3 + 16x = -8x(x^2 - 2) = -8x(x - \sqrt{2})(x + \sqrt{2}).
                \]
                Titik stasioner: $x = 0, \pm\sqrt{2}$.
                \begin{itemize}
                  \item Naik pada $(-\infty, -\sqrt{2}) \cup (0, \sqrt{2})$.
                  \item Turun pada $(-\sqrt{2}, 0) \cup (\sqrt{2}, \infty)$.
                  \item Maksimum relatif di $x = \pm\sqrt{2}$ dengan nilai $y(\pm\sqrt{2}) = -2(4) + 8(2) = 8$.
                  \item Minimum relatif di $x = 0$ dengan nilai $y(0) = 0$.
                \end{itemize}
          \item Turunan kedua:
                \[
                  y'' = -24x^2 + 16 = -8(3x^2 - 2).
                \]
                Titik belok saat $y'' = 0 \implies x = \pm\sqrt{\frac{2}{3}}$.
                \begin{itemize}
                  \item Cekung ke atas pada $\left(-\sqrt{\frac{2}{3}}, \sqrt{\frac{2}{3}}\right)$.
                  \item Cekung ke bawah pada $\left(-\infty, -\sqrt{\frac{2}{3}}\right) \cup \left(\sqrt{\frac{2}{3}}, \infty\right)$.
                \end{itemize}
        \end{enumerate}
\end{soaljawab}

\begin{soaljawab}{3}
\item Misalkan $\displaystyle F(x) = \int_{4}^{x} \frac{t+3}{t+2} \, dt$ untuk $x \in (-2, \infty)$. Dapatkan
        \begin{enumerate}
          \item $F(4)$
          \item $F'(4)$
          \item $F''(4)$
        \end{enumerate}
\jawaban
Berdasarkan Teorema Fundamental Kalkulus I:
        \begin{enumerate}
          \item $F(4) = \displaystyle\int_{4}^{4} \frac{t+3}{t+2} \, dt = 0$.
          \item $F'(x) = \dfrac{x+3}{x+2} \implies F'(4) = \dfrac{4+3}{4+2} = \dfrac{7}{6}$.
          \item $F''(x) = \dfrac{d}{dx}\left(\frac{x+3}{x+2}\right) = \dfrac{1(x+2) - (x+3)(1)}{(x+2)^2} = \dfrac{-1}{(x+2)^2}$.
                Maka $F''(4) = \dfrac{-1}{(4+2)^2} = -\dfrac{1}{36}$.
        \end{enumerate}
\end{soaljawab}

\begin{soaljawab}{4}
\item Diketahui
        \[
          z = \frac{(4 \operatorname{cis} 40^\circ)^3}{(2 \operatorname{cis} 60^\circ)^4}.
        \]
        \begin{enumerate}
          \item Sederhanakan $z$ dalam satu bentuk kutub $r\operatorname{cis}\theta$ dimana $0 \le \theta < 2\pi$.
          \item Carilah semua bilangan kompleks $w$ yang memenuhi $w = z^{1/2}$.
        \end{enumerate}
\jawaban
\begin{enumerate}
          \item Sederhanakan $z$:
                \[
                  (4 \operatorname{cis} 40^\circ)^3 = 4^3 \operatorname{cis}(3 \times 40^\circ) = 64 \operatorname{cis}(120^\circ).
                \]
                \[
                  (2 \operatorname{cis} 60^\circ)^4 = 2^4 \operatorname{cis}(4 \times 60^\circ) = 16 \operatorname{cis}(240^\circ).
                \]
                Maka:
                \[
                  z = \frac{64}{16} \operatorname{cis}(120^\circ - 240^\circ) = 4 \operatorname{cis}(-120^\circ) = 4 \operatorname{cis}(240^\circ) = 4 \operatorname{cis}\left(\frac{4\pi}{3}\right).
                \]
          \item Nilai $w = z^{1/2}$:
                \[
                  w_k = \sqrt{4} \operatorname{cis}\left(\frac{240^\circ + 360^\circ k}{2}\right) = 2 \operatorname{cis}(120^\circ + 180^\circ k) \quad \text{untuk } k = 0, 1.
                \]
                \begin{align*}
                  w_0 & = 2 \operatorname{cis}(120^\circ) = 2\left(-\frac{1}{2} + i\frac{\sqrt{3}}{2}\right) = -1 + i\sqrt{3}, \\
                  w_1 & = 2 \operatorname{cis}(300^\circ) = 2\left(\frac{1}{2} - i\frac{\sqrt{3}}{2}\right) = 1 - i\sqrt{3}.
                \end{align*}
        \end{enumerate}
\end{soaljawab}

\begin{soaljawab}{5}
\item Diberikan sistem persamaan linier berikut
        \[
          \begin{aligned}
            x_1 + x_2 + 2x_3    & = 4,  \\
            2x_1 + 3x_2 + 6x_3  & = 10, \\
            3x_1 + 6x_2 + 10x_3 & = 17.
          \end{aligned}
        \]
        \begin{enumerate}
          \item Nyatakan sistem tersebut dalam bentuk $A\mathbf{x} = \mathbf{b}$ dan dapatkan determinan dari $A$.
          \item Dapatkan $x_2$ dengan aturan Cramer.
        \end{enumerate}
\jawaban
\begin{enumerate}
\item Matriks koefisien, vektor variabel, dan ruas kanan adalah
\[
    A=\begin{bmatrix}1&1&2\\2&3&6\\3&6&10\end{bmatrix},\qquad
    \mathbf{x}=\begin{bmatrix}x_1\\x_2\\x_3\end{bmatrix},\qquad
    \mathbf{b}=\begin{bmatrix}4\\10\\17\end{bmatrix}.
\]
Dengan demikian, bentuk $A\mathbf{x}=\mathbf{b}$ ialah
\[
    \begin{bmatrix}1&1&2\\2&3&6\\3&6&10\end{bmatrix}
    \begin{bmatrix}x_1\\x_2\\x_3\end{bmatrix}
    =\begin{bmatrix}4\\10\\17\end{bmatrix}.
\]
Gunakan ekspansi kofaktor pada \textbf{baris pertama}:
\[
    \det(A)=a_{11}C_{11}+a_{12}C_{12}+a_{13}C_{13},\qquad
    C_{ij}=(-1)^{i+j}\det(M_{ij}),
\]
dengan $M_{ij}$ diperoleh dengan menghapus baris $i$ dan kolom $j$.
Tanda kofaktor pada baris pertama adalah $+,-,+$, sehingga
\begin{align*}
    \det(A)
    &=1\begin{vmatrix}3&6\\6&10\end{vmatrix}
      -1\begin{vmatrix}2&6\\3&10\end{vmatrix}
      +2\begin{vmatrix}2&3\\3&6\end{vmatrix}\\
    &=1(3\cdot10-6\cdot6)-1(2\cdot10-6\cdot3)
      +2(2\cdot6-3\cdot3)\\
    &=1(30-36)-1(20-18)+2(12-9)\\
    &=-6-2+6=-2.
\end{align*}
Karena $\det(A)\ne0$, sistem memiliki solusi tunggal dan aturan Cramer dapat digunakan.

\item Untuk mencari $x_2$, ganti \textbf{kolom kedua} dari $A$ dengan $\mathbf b$:
\[
    A_2=\begin{bmatrix}1&4&2\\2&10&6\\3&17&10\end{bmatrix}.
\]
Ekspansi kofaktor $\det(A_2)$ pada baris pertama memberikan
\begin{align*}
    \det(A_2)
    &=1\begin{vmatrix}10&6\\17&10\end{vmatrix}
      -4\begin{vmatrix}2&6\\3&10\end{vmatrix}
      +2\begin{vmatrix}2&10\\3&17\end{vmatrix}\\
    &=1(10\cdot10-6\cdot17)-4(2\cdot10-6\cdot3)
      +2(2\cdot17-10\cdot3)\\
    &=1(100-102)-4(20-18)+2(34-30)\\
    &=-2-8+8=-2.
\end{align*}
Jadi, menurut aturan Cramer,
\[
    x_2=\frac{\det(A_2)}{\det(A)}=\frac{-2}{-2}=1.
\]
\end{enumerate}
\end{soaljawab}

\label{akhir-paket-7}

% PAKET 8: Selasa, 9 Desember 2025 | 09.00 - 10.40 WIB | Mengulang (Kelas 115, 116)
\startpacket{8}
\pdfbookmark[0]{Paket 8: Selasa 09.00 Mengulang (Kelas 115, 116)}{paket8}
\examheaderID{Kalkulus 1, 3 SKS}{Selasa, 9 Desember 2025}{09:00 -- 10:40, 100 menit}{115, 116, Tertutup}
\warningsID
\cpmktableID
\scoretableID
\begin{soaljawab}{1}
\item Hitunglah integral berikut
        \[
          \int 3x^2(x^3 - 6)^7 \, dx.
        \]
\jawaban
Substitusi $u = x^3 - 6 \implies du = 3x^2\,dx$:
        \[
          \int 3x^2(x^3 - 6)^7 \, dx = \int u^7 \, du = \frac{1}{8} u^8 + C = \frac{1}{8}(x^3 - 6)^8 + C.
        \]
\end{soaljawab}

\begin{soaljawab}{2}
\item Dapatkan turunan pertama, titik stasioner, dan maksimum dan minimum relatif dari
        \[
          f(x) = x^4 - 8x^2 + 2.
        \]
\jawaban
\begin{itemize}
          \item Turunan pertama: $f'(x) = 4x^3 - 16x = 4x(x^2 - 4) = 4x(x - 2)(x + 2)$.
          \item Titik stasioner: $x = 0, \pm 2$.
          \item Turunan kedua: $f''(x) = 12x^2 - 16$.
                \begin{itemize}
                  \item $x = 0$: $f''(0) = -16 < 0 \implies$ maksimum relatif di $(0, 2)$.
                  \item $x = \pm 2$: $f''(\pm 2) = 12(4) - 16 = 32 > 0 \implies$ minimum relatif di $(\pm 2, -14)$ karena $f(\pm 2) = 16 - 32 + 2 = -14$.
                \end{itemize}
        \end{itemize}
\end{soaljawab}

\begin{soaljawab}{3}
\item Gambarkan grafik dari $y = \dfrac{1}{2}x + 2$ untuk $2 \le x \le 6$ dan hitung integral
        \[
          \int_{2}^{6} \left(\frac{1}{2}x + 2\right) \, dx
        \]
        dengan menggunakan rumus dari geometri bidang yang sesuai.
\jawaban
Daerah di bawah kurva $y = \frac{1}{2}x + 2$ pada interval $[2, 6]$ adalah trapesium dengan:
        \begin{itemize}
          \item Sisi sejajar: $y(2) = 1 + 2 = 3$ dan $y(6) = 3 + 2 = 5$
          \item Tinggi trapesium: $t = 6 - 2 = 4$
        \end{itemize}
        Luas trapesium:
        \[
          \int_{2}^{6} \left(\frac{1}{2}x + 2\right) \, dx = \frac{3 + 5}{2} \cdot 4 = 16.
        \]
\end{soaljawab}

\begin{soaljawab}{4}
\item Sederhanakan
        \[
          \frac{5 + 5i}{3 - 4i} + \frac{20i}{4 + 3i}
        \]
        dalam bentuk $a + bi$.
\jawaban
Hitung masing-masing pecahan:
        \[
          \frac{5 + 5i}{3 - 4i} = \frac{(5 + 5i)(3 + 4i)}{3^2 + (-4)^2} = \frac{15 + 20i + 15i - 20}{25} = \frac{-5 + 35i}{25} = -\frac{1}{5} + \frac{7}{5}i.
        \]
        \[
          \frac{20i}{4 + 3i} = \frac{20i(4 - 3i)}{4^2 + 3^2} = \frac{80i - 60i^2}{25} = \frac{60 + 80i}{25} = \frac{12}{5} + \frac{16}{5}i.
        \]
        Jumlahkan:
        \[
          \left(-\frac{1}{5} + \frac{7}{5}i\right) + \left(\frac{12}{5} + \frac{16}{5}i\right) = \frac{11}{5} + \frac{23}{5}i.
        \]
\end{soaljawab}

\begin{soaljawab}{5}
\item Diketahui $A = \begin{pmatrix} 1 & 4 \\ 4 & 15 \end{pmatrix}$.
        \begin{enumerate}
          \item Dapatkan $A^{-1}$ dengan menggunakan rumus $\dfrac{1}{ad-bc} \begin{pmatrix} d & -b \\ -c & a \end{pmatrix}$.
          \item Dapatkan $A^{-1}$ dengan menggunakan operasi baris elementer.
        \end{enumerate}
\jawaban
\begin{enumerate}
\item Untuk $A=\begin{pmatrix}a&b\\c&d\end{pmatrix}$, rumus inversnya adalah
\[
    A^{-1}=\frac1{ad-bc}\begin{pmatrix}d&-b\\-c&a\end{pmatrix},
    \qquad ad-bc\ne0.
\]
Ekspansi kofaktor pada baris pertama memberi
\[
    \det(A)=\begin{vmatrix}
            1 & 4 \\
            4 & 15
        \end{vmatrix}=1\cdot15-4\cdot4=15-16=-1.
\]
Karena determinan tidak nol, inversnya ada. Substitusi ke rumus memberikan
\[
    A^{-1}=\frac1{-1}\begin{pmatrix}
            15 & -4 \\
            -4 & 1
        \end{pmatrix}=\begin{pmatrix}
            -15 & 4 \\
            4 & -1
        \end{pmatrix}.
\]
\item Untuk metode OBE, ubah $[A\mid I_2]$ menjadi $[I_2\mid A^{-1}]$.
Notasi $B_i$ menyatakan baris ke-$i$.
{\setlength{\arraycolsep}{4bp}
\begin{align*}
    & \left[\begin{array}{cc|cc}
            1 & 4 & 1 & 0 \\
            4 & 15 & 0 & 1
        \end{array}\right]
        \begin{array}{c}\xrightarrow{B_2-4B_1}\end{array}
        \left[\begin{array}{cc|cc}
            1 & 4 & 1 & 0 \\
            0 & -1 & -4 & 1
        \end{array}\right] \\[6bp]
    & \begin{array}{c}\xrightarrow{-B_2}\end{array}
        \left[\begin{array}{cc|cc}
            1 & 4 & 1 & 0 \\
            0 & 1 & 4 & -1
        \end{array}\right]
        \begin{array}{c}\xrightarrow{B_1-4B_2}\end{array}
        \left[\begin{array}{cc|cc}
            1 & 0 & -15 & 4 \\
            0 & 1 & 4 & -1
        \end{array}\right]
\end{align*}}
Bagian kiri telah menjadi $I_2$. Maka bagian kanan memberi hasil yang sama:
\[
    A^{-1}=\begin{pmatrix}
            -15 & 4 \\
            4 & -1
        \end{pmatrix}.
\]
\end{enumerate}
\end{soaljawab}

\label{akhir-paket-8}

% PAKET 9: Selasa, 9 Desember 2025 | 11.00 - 12.40 WIB | Mengulang (Kelas 102, 108-112)
\startpacket{9}
\pdfbookmark[0]{Paket 9: Selasa 11.00 Mengulang (Kelas 102, 108--112)}{paket9}
\examheaderID{Kalkulus 1, 3 SKS}{Selasa, 9 Desember 2025}{11:00 -- 12:40, 100 menit}{102, 108 -- 112, Tertutup}
\warningsID
\cpmktableID
\scoretableID
\begin{soaljawab}{1}
\item Hitunglah integral berikut
        \[
          \int 4x^3(x^4 - 2)^6 \, dx.
        \]
\jawaban
Substitusi $u = x^4 - 2 \implies du = 4x^3\,dx$:
        \[
          \int 4x^3(x^4 - 2)^6 \, dx = \int u^6 \, du = \frac{1}{7} u^7 + C = \frac{1}{7}(x^4 - 2)^7 + C.
        \]
\end{soaljawab}

\begin{soaljawab}{2}
\item Diketahui $f(x) = 5 + 12x - x^3$.
        \begin{enumerate}
          \item Dapatkan turunan pertama, titik stasioner, dan selang naik turun dari fungsi tersebut.
          \item Dapatkan turunan kedua, titik belok, dan kecekungan dari fungsi tersebut.
        \end{enumerate}
\jawaban
\begin{enumerate}
          \item Turunan pertama:
                \[
                  f'(x) = 12 - 3x^2 = 3(4 - x^2) = 3(2 - x)(2 + x).
                \]
                Titik stasioner: $x = \pm 2$.
                \begin{itemize}
                  \item Naik pada interval $(-2, 2)$.
                  \item Turun pada interval $(-\infty, -2) \cup (2, \infty)$.
                \end{itemize}
          \item Turunan kedua:
                \[
                  f''(x) = -6x.
                \]
                Titik belok di $x = 0$ dengan koordinat $(0, f(0)) = (0, 5)$.
                \begin{itemize}
                  \item Cekung ke atas pada $(-\infty, 0)$ (karena $f''(x) > 0$).
                  \item Cekung ke bawah pada $(0, \infty)$ (karena $f''(x) < 0$).
                \end{itemize}
        \end{enumerate}
\end{soaljawab}

\begin{soaljawab}{3}
\item Gambarkan grafik dari $y = -2x + 8$ untuk $1 \le x \le 3$ dan hitung integral
        \[
          \int_{1}^{3} (-2x + 8) \, dx
        \]
        dengan menggunakan rumus dari geometri bidang yang sesuai.
\jawaban
Daerah di bawah garis $y = -2x + 8$ pada interval $[1, 3]$ merupakan trapesium:
        \begin{itemize}
          \item Sisi sejajar: $y(1) = -2 + 8 = 6$ dan $y(3) = -6 + 8 = 2$
          \item Tinggi trapesium: $t = 3 - 1 = 2$
        \end{itemize}
        Luas:
        \[
          \int_{1}^{3} (-2x + 8) \, dx = \frac{6 + 2}{2} \cdot 2 = 8.
        \]
\end{soaljawab}

\begin{soaljawab}{4}
\item Sederhanakan
        \[
          \frac{10 + 5i}{4 - 3i} + \frac{15i}{3 + 4i}
        \]
        dalam bentuk $a + bi$.
\jawaban
Sederhanakan masing-masing:
        \[
          \frac{10 + 5i}{4 - 3i} = \frac{(10 + 5i)(4 + 3i)}{4^2 + (-3)^2} = \frac{40 + 30i + 20i - 15}{25} = \frac{25 + 50i}{25} = 1 + 2i.
        \]
        \[
          \frac{15i}{3 + 4i} = \frac{15i(3 - 4i)}{3^2 + 4^2} = \frac{45i - 60i^2}{25} = \frac{60 + 45i}{25} = \frac{12}{5} + \frac{9}{5}i.
        \]
        Jumlahkan:
        \[
          (1 + 2i) + \left(\frac{12}{5} + \frac{9}{5}i\right) = \frac{17}{5} + \frac{19}{5}i.
        \]
\end{soaljawab}

\begin{soaljawab}{5}
\item Diketahui $A = \begin{pmatrix} 1 & 5 \\ 3 & 14 \end{pmatrix}$.
        \begin{enumerate}
          \item Dapatkan $A^{-1}$ dengan menggunakan rumus $\dfrac{1}{ad-bc} \begin{pmatrix} d & -b \\ -c & a \end{pmatrix}$.
          \item Dapatkan $A^{-1}$ dengan menggunakan operasi baris elementer.
        \end{enumerate}
\jawaban
\begin{enumerate}
\item Untuk $A=\begin{pmatrix}a&b\\c&d\end{pmatrix}$, rumus inversnya adalah
\[
    A^{-1}=\frac1{ad-bc}\begin{pmatrix}d&-b\\-c&a\end{pmatrix},
    \qquad ad-bc\ne0.
\]
Ekspansi kofaktor pada baris pertama memberi
\[
    \det(A)=\begin{vmatrix}
            1 & 5 \\
            3 & 14
        \end{vmatrix}=1\cdot14-5\cdot3=14-15=-1.
\]
Karena determinan tidak nol, inversnya ada. Substitusi ke rumus memberikan
\[
    A^{-1}=\frac1{-1}\begin{pmatrix}
            14 & -5 \\
            -3 & 1
        \end{pmatrix}=\begin{pmatrix}
            -14 & 5 \\
            3 & -1
        \end{pmatrix}.
\]
\item Untuk metode OBE, ubah $[A\mid I_2]$ menjadi $[I_2\mid A^{-1}]$.
Notasi $B_i$ menyatakan baris ke-$i$.
{\setlength{\arraycolsep}{4bp}
\begin{align*}
    & \left[\begin{array}{cc|cc}
            1 & 5 & 1 & 0 \\
            3 & 14 & 0 & 1
        \end{array}\right]
        \begin{array}{c}\xrightarrow{B_2-3B_1}\end{array}
        \left[\begin{array}{cc|cc}
            1 & 5 & 1 & 0 \\
            0 & -1 & -3 & 1
        \end{array}\right] \\[6bp]
    & \begin{array}{c}\xrightarrow{-B_2}\end{array}
        \left[\begin{array}{cc|cc}
            1 & 5 & 1 & 0 \\
            0 & 1 & 3 & -1
        \end{array}\right]
        \begin{array}{c}\xrightarrow{B_1-5B_2}\end{array}
        \left[\begin{array}{cc|cc}
            1 & 0 & -14 & 5 \\
            0 & 1 & 3 & -1
        \end{array}\right]
\end{align*}}
Bagian kiri telah menjadi $I_2$. Maka bagian kanan memberi hasil yang sama:
\[
    A^{-1}=\begin{pmatrix}
            -14 & 5 \\
            3 & -1
        \end{pmatrix}.
\]
\end{enumerate}
\end{soaljawab}

\label{akhir-paket-9}

% PAKET 10: Selasa, 9 Desember 2025 | 15.30 - 17.10 WIB | Reguler (Kelas 30-39)
\startpacket{10}
\pdfbookmark[0]{Paket 10: Selasa 15.30 Reguler (Kelas 30--39)}{paket10}
\examheaderID{Kalkulus 1, 3 SKS}{Selasa, 9 Desember 2025}{15:30 -- 17:10, 100 menit}{30 -- 39, Tertutup}
\warningsID
\cpmktableID
\scoretableID
\begin{soaljawab}{1}
\item Hitunglah integral berikut
        \[
          \int x^3 \sqrt{x^2 + 5} \, dx.
        \]
\jawaban
Tuliskan $x^3 = x^2 \cdot x$. Substitusi $u = x^2 + 5 \implies x^2 = u - 5$, dan $du = 2x\,dx \implies x\,dx = \frac{1}{2}\,du$.
        \begin{align*}
          \int x^3 \sqrt{x^2 + 5} \, dx & = \int (u - 5)\sqrt{u} \left(\frac{1}{2}\,du\right) = \frac{1}{2} \int (u^{3/2} - 5u^{1/2}) \, du \\
                                        & = \frac{1}{2} \left(\frac{2}{5}u^{5/2} - 5 \cdot \frac{2}{3}u^{3/2}\right) + C                    \\
                                        & = \frac{1}{5}(x^2 + 5)^{5/2} - \frac{5}{3}(x^2 + 5)^{3/2} + C.
        \end{align*}
\end{soaljawab}

\begin{soaljawab}{2}
\item Tentukan nilai maksimum dan minimum absolut dari
        \[
          f(x) = \begin{cases}
            4x - 2,       & x < 1,   \\
            x^2 - 5x + 6, & x \ge 1,
          \end{cases}
        \]
        pada interval $\left[\frac{1}{2}, 3\right]$.
\jawaban
Evaluasi fungsi pada domain $\left[\frac{1}{2}, 3\right]$:
        \begin{itemize}
          \item Periksa kekontinuan di titik batas cabang $x = 1$:
                $\displaystyle \lim_{x \to 1^-} (4x - 2) = 2$, dan $f(1) = 1^2 - 5(1) + 6 = 2$. Jadi $f(x)$ kontinu di $x = 1$.
          \item Pada $\left[\frac{1}{2}, 1\right)$: $f'(x) = 4 \ne 0$ (tidak ada titik stasioner).
          \item Pada $(1, 3]$: $f'(x) = 2x - 5$. Titik stasioner $f'(x) = 0 \implies x = \frac{5}{2} = 2.5 \in (1, 3)$.
          \item Evaluasi nilai-nilai calon ekstrim:
                \begin{itemize}
                  \item $x = \frac{1}{2}$: $f\left(\frac{1}{2}\right) = 4\left(\frac{1}{2}\right) - 2 = 0$.
                  \item $x = 1$: $f(1) = 2$.
                  \item $x = \frac{5}{2}$: $f\left(\frac{5}{2}\right) = \left(\frac{5}{2}\right)^2 - 5\left(\frac{5}{2}\right) + 6 = \frac{25}{4} - \frac{25}{2} + 6 = -\frac{1}{4} = -0.25$.
                  \item $x = 3$: $f(3) = 3^2 - 5(3) + 6 = 0$.
                \end{itemize}
        \end{itemize}
        Dari nilai-nilai di atas:
        \begin{itemize}
          \item Nilai maksimum absolut adalah $2$ pada $x = 1$.
          \item Nilai minimum absolut adalah $-\dfrac{1}{4}$ pada $x = \dfrac{5}{2}$.
        \end{itemize}
\end{soaljawab}

\begin{soaljawab}{3}
\item Misalkan $\displaystyle F(x) = \int_{4}^{x} \sqrt{t^2 + 9} \, dt$. Dapatkan
        \begin{enumerate}
          \item $F(4)$
          \item $F'(4)$
          \item $F''(4)$
        \end{enumerate}
\jawaban
Berdasarkan Teorema Fundamental Kalkulus I:
        \begin{enumerate}
          \item $F(4) = \displaystyle\int_{4}^{4} \sqrt{t^2 + 9} \, dt = 0$.
          \item $F'(x) = \sqrt{x^2 + 9} \implies F'(4) = \sqrt{4^2 + 9} = \sqrt{25} = 5$.
          \item $F''(x) = \dfrac{d}{dx}\sqrt{x^2 + 9} = \dfrac{2x}{2\sqrt{x^2 + 9}} = \dfrac{x}{\sqrt{x^2 + 9}}$.
                Maka $F''(4) = \dfrac{4}{\sqrt{4^2 + 9}} = \dfrac{4}{5}$.
        \end{enumerate}
\end{soaljawab}

\begin{soaljawab}{4}
\item Carilah semua bilangan kompleks $z$ yang memenuhi $z = (-16)^{1/4}$. Gambarlah semua bilangan kompleks yang memenuhi tersebut pada bidang kompleks.
\jawaban
Nyatakan $-16$ dalam bentuk polar:
        \[
          -16 = 16 \operatorname{cis}(180^\circ) = 16 e^{i\pi}.
        \]
        Akar keempatnya untuk $k = 0, 1, 2, 3$:
        \[
          z_k = \sqrt[4]{16} \operatorname{cis}\left(\frac{180^\circ + 360^\circ k}{4}\right) = 2 \operatorname{cis}(45^\circ + 90^\circ k).
        \]
        Diperoleh ke-4 akar:
        \begin{align*}
          z_0 & = 2 \operatorname{cis}(45^\circ) = 2\left(\frac{\sqrt{2}}{2} + i\frac{\sqrt{2}}{2}\right) = \sqrt{2} + i\sqrt{2},    \\
          z_1 & = 2 \operatorname{cis}(135^\circ) = 2\left(-\frac{\sqrt{2}}{2} + i\frac{\sqrt{2}}{2}\right) = -\sqrt{2} + i\sqrt{2}, \\
          z_2 & = 2 \operatorname{cis}(225^\circ) = 2\left(-\frac{\sqrt{2}}{2} - i\frac{\sqrt{2}}{2}\right) = -\sqrt{2} - i\sqrt{2}, \\
          z_3 & = 2 \operatorname{cis}(315^\circ) = 2\left(\frac{\sqrt{2}}{2} - i\frac{\sqrt{2}}{2}\right) = \sqrt{2} - i\sqrt{2}.
        \end{align*}
        Pada bidang kompleks, keempat titik tersebut terletak pada lingkaran berpusat di titik asal dengan jari-jari $r = 2$ dan membentuk persegi yang simetris di empat kuadran.
\end{soaljawab}

\begin{soaljawab}{5}
\item Hitunglah determinan dari
        \[
          A = \begin{bmatrix}
            1 & 2 & 1  & 2 \\
            1 & 3 & -1 & 1 \\
            0 & 2 & 2  & 3 \\
            2 & 4 & -3 & 2
          \end{bmatrix}.
        \]
\jawaban
\textbf{Langkah 1: buat nol pada kolom pertama dengan OBE.}
Operasi $B_i\leftarrow B_i+kB_j$ tidak mengubah determinan.
Lakukan $B_2\leftarrow B_2-B_1$ dan $B_4\leftarrow B_4-2B_1$:
{\setlength{\arraycolsep}{4bp}
\begin{align*}
    & \left[\begin{array}{cccc}
            1 & 2 & 1 & 2 \\
            1 & 3 & -1 & 1 \\
            0 & 2 & 2 & 3 \\
            2 & 4 & -3 & 2
        \end{array}\right]
        \begin{array}{c}\xrightarrow{B_2-B_1}\end{array}
        \left[\begin{array}{cccc}
            1 & 2 & 1 & 2 \\
            0 & 1 & -2 & -1 \\
            0 & 2 & 2 & 3 \\
            2 & 4 & -3 & 2
        \end{array}\right] \\[6bp]
    & \begin{array}{c}\xrightarrow{B_4-2B_1}\end{array}
        \left[\begin{array}{cccc}
            1 & 2 & 1 & 2 \\
            0 & 1 & -2 & -1 \\
            0 & 2 & 2 & 3 \\
            0 & 0 & -5 & -2
        \end{array}\right]
\end{align*}}
Sebut matriks terakhir $U$. Karena kedua operasi hanya menambahkan kelipatan
baris lain, $\det(A)=\det(U)$.

\smallskip
\textbf{Langkah 2: ekspansi kofaktor pada kolom pertama $U$.}
Dengan $C_{ij}=(-1)^{i+j}\det(M_{ij})$, diperoleh
\begin{align*}
    \det(A)=\det(U)
    &=u_{11}C_{11}+u_{21}C_{21}+u_{31}C_{31}+u_{41}C_{41}\\
    &=1C_{11}+0C_{21}+0C_{31}+0C_{41}\\
    &=1(-1)^{1+1}\begin{vmatrix}
        1&-2&-1\\2&2&3\\0&-5&-2
      \end{vmatrix}.
\end{align*}
Minor di atas diperoleh dengan menghapus baris pertama dan kolom pertama $U$.

\smallskip
\textbf{Langkah 3: ekspansi determinan $3\times3$ pada baris pertama.}
Tanda kofaktor pada baris pertama adalah $+,-,+$, sehingga
\begin{align*}
    \det(A)
    &=1\begin{vmatrix}2&3\\-5&-2\end{vmatrix}
      -(-2)\begin{vmatrix}2&3\\0&-2\end{vmatrix}
      +(-1)\begin{vmatrix}2&2\\0&-5\end{vmatrix}\\
    &=1\bigl(2(-2)-3(-5)\bigr)
      +2\bigl(2(-2)-3\cdot0\bigr)
      -\bigl(2(-5)-2\cdot0\bigr)\\
    &=1(-4+15)+2(-4-0)-(-10-0)\\
    &=11-8+10=13.
\end{align*}
Jadi $\det(A)=13$.
\end{soaljawab}

\label{akhir-paket-10}
\end{document}
